\chapter{Blepharospasm}
\index{Blepharospasm}
\label{sec:Blepharospasm}

\section{Overview}
\begin{itemize}[noitemsep]
  \item Blepharospasm is an involuntary, forceful, and often progressive closing of the eyelids caused by abnormal contraction of the eyelid muscles
  \item It is a form of focal dystonia and can range from occasional increased blinking to severe, functionally disabling spasms
  \item It typically begins in middle age and affects women more than men
\end{itemize}
\section{Causes}
The cause of primary (essential) blepharospasm is unknown, but it is thought to involve abnormal signal processing in the basal ganglia of the brain. It can be worsened by bright light, stress, fatigue, and eye strain. Secondary blepharospasm can follow eye disease, dry eyes, certain medications, stroke, or other neurological disorders.
\section{Risk Factors}
\begin{itemize}[noitemsep]
  \item Female sex
  \item Age between 40 and 60 years
  \item Family history of dystonia
  \item Eye irritation, dry eyes, or light sensitivity
  \item Stress and fatigue
  \item Neurological conditions (Parkinson disease, stroke, multiple sclerosis)
\end{itemize}
\section{Signs and Symptoms}
\subsection{Common symptoms}
\begin{itemize}[noitemsep]
  \item Increased frequency of blinking
  \item Involuntary, forceful eyelid closure that is difficult to stop
  \item Sensitivity to bright light (photophobia)
  \item Dry or irritated eyes
  \item Spasms that worsen with stress, fatigue, or concentration
  \item In severe cases, functional blindness because the eyes cannot be opened
\end{itemize}
\subsection{Emergency warning signs}
\begin{itemize}[noitemsep]
  \item Sudden eyelid drooping, double vision, or facial weakness may indicate a more serious neurological problem and requires evaluation
\end{itemize}
\section{Progression and Stages}
Blepharospasm often starts with increased blinking and eye irritation and progresses over months to years to more forceful and frequent spasms. It may spread to involve other facial muscles. The severity fluctuates and can be triggered by specific activities. It is not progressive to a life-threatening degree, but it can become functionally disabling.
\section{Complications}
\begin{itemize}[noitemsep]
  \item Functional visual impairment (inability to keep eyes open)
  \item Difficulty driving, reading, and working
  \item Social embarrassment and anxiety
  \item Depression and reduced quality of life
  \item Injury from falls during spasm episodes
  \item Side effects of treatment
\end{itemize}
\section{Diagnosis}
\begin{itemize}[noitemsep]
  \item Clinical diagnosis based on history and observation of the spasms
  \item Neurological examination to rule out other movement disorders
  \item Eye examination to exclude eye surface disease or ocular irritation
  \item Imaging (MRI) in atypical or secondary cases
  \item Electromyography (EMG) may support the diagnosis
\end{itemize}
\section{Prevention}
\begin{itemize}[noitemsep]
  \item Primary blepharospasm is not preventable
  \item Managing eye irritation and dry eyes
  \item Reducing stress and fatigue
  \item Using tinted lenses for light sensitivity
  \item Early recognition and referral to a neurologist
\end{itemize}
\section{Treatment Options}
\begin{itemize}[noitemsep]
  \item Botulinum toxin injections into the eyelid muscles are the first-line and most effective treatment
  \item Artificial tears and lubricants for associated dry eyes
  \item Photochromic or FL-41 tinted lenses for light sensitivity
  \item Oral medications (benzodiazepines, anticholinergics) in some cases with limited benefit
  \item Biofeedback, relaxation, and sensory "tricks" (e.g., touching the face, humming) to ease spasms
  \item Surgical myectomy (removal of eyelid muscles) in severe, refractory cases
\end{itemize}
\section{Living with It}
\begin{itemize}[noitemsep]
  \item Schedule botulinum toxin injections regularly as recommended
  \item Use tinted lenses and sunglasses outdoors
  \item Manage stress through relaxation and adequate sleep
  \item Seek support from patient organizations
  \item Discuss functional accommodations with employers
\end{itemize}
\section{Outlook}
Blepharospasm is not life-threatening, and botulinum toxin injections provide effective relief for the vast majority of patients, with repeat treatments every 3--4 months. Although there is no cure, most people maintain good function and quality of life with treatment.
